% =================================================================
% 第3.1-3.3节：软件测试环境（共享模板）
% 说明：
%   - 3.1 测试现场
%   - 3.2 测试环境概述（包含3.2.1-3.2.3基础子节）
%   - 3.2.1 软件项
%   - 3.2.2 硬件和固件项
%   - 3.2.3 其他材料
%
% 注意：后续子节（3.2.4-3.2.10）由各文档自行定义
% =================================================================

\subsection*{3.1 测试现场}
\subsection*{3.2 测试环境概述}

{\normalsize
测试环境总体架构如下：

（1）硬件环境：包括服务器、客户端、网络设备等；

（2）软件环境：包括操作系统、数据库系统、中间件等；

（3）网络环境：局域网，带宽xxxMbps；

（4）测试工具：xxxx。
}

\subsubsection*{3.2.1 软件项}

{\normalsize
测试环境中的软件项如表\ref{\GjbEnvSoftwareTableLabel}所示。

{\settablespacing
\begin{longtblr}[theme=gjb,caption={测试环境被测软件},label={\GjbEnvSoftwareTableLabel}]{
  colspec={|p{0.8cm}|p{3.5cm}|p{2.2cm}|p{2.8cm}|p{2.5cm}|p{2.7cm}|},
  rowhead=1,
  hlines={wd=\GjbTableRuleWd,fg=\GjbTableRuleColor},
  column{1}={halign=c},
}
序号 & 软件名称 & 版本号 & 用途 & 来源 & 存储介质 \\
\Seq & xxxxx & xxxx & 被测软件 & 本单位 & 硬盘 \\
\Seq & 火狐浏览器 & 123.0 & 操作页面 & 开源 & 硬盘 \\
\end{longtblr}
}
\vspace{-6pt}

测试环境中的支撑软件包括操作系统、数据库系统等，为被测软件提供必要的运行环境。
}

\subsubsection*{3.2.2 硬件和固件项}

{\normalsize
测试环境中的硬件和固件项如表\ref{\GjbEnvHardwareTableLabel}所示。

{\settablespacing
\begin{longtblr}[theme=gjb,caption={测试环境硬件项},label={\GjbEnvHardwareTableLabel}]{
  colspec={|p{0.8cm}|p{1.5cm}|p{1.8cm}|p{2.0cm}|p{0.8cm}|p{2.65cm}|p{1.0cm}|p{2.4cm}|p{1.0cm}|},
  rowhead=1,
  hlines={wd=\GjbTableRuleWd,fg=\GjbTableRuleColor},
  column{1}={halign=c},
}
序号 & 硬件名称 & 型号 & 编号/序列号 & 数量 & 配置 & 用途 & 软件部署环境 & 来源 \\
\Seq & 服务器 & 长城 & / & 1 & CPU：FT-2000+ *2 内存：256G & 测试机 & 操作系统：Kylin V10，docker & 本单位 \\
\Seq & 台式机 & 联想Think station & / & 2 & CPU：Intel Corei7-12700 内存：32G & 测试机 & 操作系统：Ubuntu 22.04，docker & 本单位 \\
\end{longtblr}
}
\vspace{-6pt}

硬件配置满足测试要求，所有设备均已安装调试完毕，运行正常。固件版本符合系统要求。
}

\subsubsection*{3.2.3 其他材料}

{\normalsize
测试所需的其他材料包括：

（1）测试数据：xxxx；

（2）测试用例：xxxx；

（3）用户手册：xxxx；

（4）技术文档：xxxx。
}
